\documentclass[12pt,a4paper]{article}

\usepackage[T2A]{fontenc}
\usepackage[utf8]{inputenc}
\usepackage[russian]{babel}
\usepackage{amsmath}
\usepackage{amssymb}
\usepackage{geometry}
\usepackage{hyperref}
\usepackage{xcolor}
\usepackage{listings}
\usepackage{booktabs}

\geometry{left=25mm,right=20mm,top=20mm,bottom=25mm}

\hypersetup{
  colorlinks=true,
  linkcolor=blue,
  urlcolor=blue
}

\lstdefinelanguage{JavaScript}{
  keywords={
    const,let,var,function,return,if,else,for,while,do,import,from,export,
    default,new,class,extends,throw,true,false,null,undefined,Math,Array,Number
  },
  sensitive=true,
  comment=[l]{//},
  morecomment=[s]{/*}{*/},
  morestring=[b]",
  morestring=[b]',
  morestring=[b]`
}

\lstdefinestyle{jsstyle}{
  basicstyle=\ttfamily\small,
  breaklines=true,
  frame=single,
  columns=fullflexible,
  keepspaces=true,
  showstringspaces=false,
  keywordstyle=\color{blue!60!black},
  commentstyle=\color{green!40!black},
  stringstyle=\color{red!50!black}
}

\lstset{style=jsstyle}

\title{
  \textbf{Отчет по вычислительному практикуму}\\
  \large Моделирование колец Ньютона для линзы заданного радиуса
}
\author{Студент: \underline{Лазимов Даниил M3203}}
\date{\today}

\begin{document}

\maketitle

\tableofcontents
\newpage

\section{Цель работы}

Цель работы состоит в построении компьютерной модели колец Ньютона для тонкой пленки между
сферической линзой и плоской пластиной. В модели требуется получить цветное распределение
интенсивности интерференционной картины, построить график зависимости интенсивности от радиальной
координаты и исследовать, как картина меняется при переходе от монохроматического света к
квазимонохроматическому.

Отдельное внимание уделяется тому, чтобы модель была не только визуально наглядной, но и физически
проверяемой: радиусы колец должны согласовываться с аналитическими формулами, а конечная ширина
спектра должна приводить к ожидаемому уменьшению контраста дальних колец.

\section{Постановка задачи}

Рассматривается классическая схема наблюдения колец Ньютона: сферическая поверхность линзы лежит на
плоской стеклянной пластине. Между ними образуется тонкая пленка, толщина которой зависит от
расстояния до центра контакта. Интерференция лучей, отраженных от разных границ пленки, дает систему
концентрических темных и светлых колец.

В программе пользователь задает:
\begin{itemize}
  \item радиус кривизны линзы \(R\);
  \item радиус наблюдаемой области;
  \item центральную длину волны \(\lambda_0\);
  \item ширину спектра \(\Delta\lambda\) в нанометрах;
  \item начальный зазор \(t_0\);
  \item показатель преломления пленки \(n\);
  \item видность интерференции \(V\);
  \item режим наблюдения: отраженный или проходящий свет.
\end{itemize}

Результатом работы программы являются две связанные визуализации: двумерная цветная картина колец и
график \(I(r)\), где \(r\) -- радиальная координата.

\section{Физическая и математическая модель}

Выбранная модель основана на двухлучевой интерференции в тонкой пленке. Такой подход является
естественным для колец Ньютона, поскольку основная наблюдаемая картина определяется разностью фаз
между двумя лучами: один отражается от верхней границы пленки, другой -- от нижней. Более сложные
эффекты, например многократные отражения, поляризация или дисперсия стекла, в данной работе не
учитываются. Для учебного моделирования это разумное приближение: оно дает стандартные формулы
радиусов колец и правильно описывает темный центр в отраженном свете.

Толщина пленки определяется геометрией сферической поверхности. Пусть \(R\) -- радиус кривизны линзы,
\(r\) -- расстояние от центра картины, \(t_0\) -- возможный начальный зазор. Тогда точная толщина
пленки равна
\[
t(r)=t_0+R-\sqrt{R^2-r^2}.
\]

При \(r\ll R\) эта формула переходит в привычное параксиальное приближение:
\[
t(r)\approx t_0+\frac{r^2}{2R}.
\]

Оптическая разность хода для почти нормального падения равна
\[
\delta(r)=2nt(r),
\]
а соответствующая фаза для длины волны \(\lambda\):
\[
\varphi(\lambda,r)=\frac{4\pi nt(r)}{\lambda}.
\]

В отраженном свете один из лучей получает дополнительный скачок фазы на \(\pi\). Поэтому при
нулевом зазоре центр картины оказывается темным. Нормированная интенсивность записывается в виде
\[
I_{\mathrm{ref}}(\lambda,r)=\frac{1-V\cos\varphi(\lambda,r)}{2}.
\]

В проходящем свете картина комплементарна:
\[
I_{\mathrm{tr}}(\lambda,r)=\frac{1+V\cos\varphi(\lambda,r)}{2}.
\]

Здесь \(V\) -- видность интерференции. При \(V=1\) получается идеальная интерференционная картина, а
при \(V=0\) кольца исчезают и остается равномерная интенсивность.

\section{Основные формулы и обозначения}

Темные кольца в отраженном свете задаются условием
\[
2nt(r_m)=m\lambda,\qquad m=0,1,2,\ldots
\]
а светлые кольца:
\[
2nt(r_m)=\left(m+\frac{1}{2}\right)\lambda.
\]

Если \(t_0=0\) и \(r\ll R\), то для темных колец получается простая формула
\[
r_m^2\approx \frac{m\lambda R}{n}.
\]

Для диаметра \(D_m=2r_m\):
\[
D_m^2\approx \frac{4m\lambda R}{n}.
\]

Именно это соотношение обычно используется в лабораторной обработке экспериментальных данных:
строится зависимость \(D_m^2\) от номера кольца \(m\), и по ее наклону можно найти длину волны или
радиус кривизны линзы.

Для квазимонохроматического источника используется гауссов спектр. Пользователь задает середину
спектра \(\lambda_0\) и ширину \(\Delta\lambda\), понимаемую как FWHM:
\[
S(\lambda)=\exp\left(-\frac{(\lambda-\lambda_0)^2}{2\sigma^2}\right),
\]
\[
\sigma=\frac{\Delta\lambda}{2\sqrt{2\ln2}}.
\]

Итоговая интенсивность получается спектральным усреднением:
\[
I(r)=
\frac{\int S(\lambda)I(\lambda,r)\,d\lambda}
{\int S(\lambda)\,d\lambda}.
\]

Для оценки того, насколько быстро исчезает контраст дальних колец, удобно использовать порядок
величины длины когерентности:
\[
L_c\approx \frac{\lambda_0^2}{\Delta\lambda}.
\]

\section{Численный метод}

Основная численная операция в работе -- дискретное спектральное усреднение. Для монохроматического
света спектр состоит из одного отсчета с весом 1. Для квазимонохроматического света строится сетка
длин волн в видимом диапазоне \(380\)--\(780\) нм, а каждому отсчету присваивается гауссов вес.
После нормировки сумма весов равна единице.

Для каждого радиуса \(r\) вычисляются:
\begin{enumerate}
  \item толщина пленки \(t(r)\);
  \item фаза \(\varphi(\lambda,r)\) для каждой длины волны;
  \item монохроматическая интенсивность \(I(\lambda,r)\);
  \item взвешенная сумма интенсивностей по спектру.
\end{enumerate}

Цветовая визуализация строится похожим способом, но вместо одной скалярной интенсивности считается
спектральная сумма в пространстве \(XYZ\):
\[
XYZ(r)=\int S(\lambda)I(\lambda,r)
\begin{bmatrix}
\bar{x}(\lambda)\\
\bar{y}(\lambda)\\
\bar{z}(\lambda)
\end{bmatrix}
d\lambda.
\]
Затем цвет переводится в линейный \(sRGB\) и гамма-корректируется для вывода на Canvas.

Так как картина осесимметрична, интенсивность зависит только от \(r\), а не от угла. Поэтому в
визуализации не требуется заново считать спектральный интеграл для каждого пикселя. Программа
сначала строит таблицу значений по радиусу, а затем заполняет пиксели по готовой таблице. Это
ускоряет интерактивную работу, но не меняет физическую модель.

\section{Алгоритм программы}

Алгоритм работы демо можно описать следующим образом:
\begin{enumerate}
  \item Считать параметры, заданные пользователем.
  \item Определить режим наблюдения: отраженный или проходящий свет.
  \item Построить спектр: один отсчет для монохроматического света или 81 отсчет для
  квазимонохроматического.
  \item Для радиальной сетки вычислить толщину пленки \(t(r)\).
  \item Для каждой длины волны из спектра вычислить фазу и интенсивность.
  \item Просуммировать вклады с учетом спектральных весов.
  \item Перевести спектральное распределение в цвет \(sRGB\).
  \item Нарисовать двумерную картину колец на Canvas.
  \item Построить график \(I(r)\) для выбранных параметров.
  \item Обновить числовые подписи: тип света, интенсивность в центре, радиус первого темного кольца,
  оценку длины когерентности.
\end{enumerate}

\section{Реализация}

Проект написан на HTML, CSS и JavaScript без внешних npm-зависимостей. Физическая часть отделена от
отрисовки, поэтому основные формулы можно тестировать независимо от интерфейса.

\begin{table}[h]
\centering
\begin{tabular}{ll}
\toprule
Файл & Назначение \\
\midrule
\texttt{index.html} & структура интерфейса \\
\texttt{src/styles.css} & визуальное оформление \\
\texttt{src/physics.js} & физическая модель, спектр, интенсивность, цвет \\
\texttt{src/app.js} & чтение параметров и Canvas-визуализация \\
\texttt{tests/physics.test.js} & автоматические проверки физических соотношений \\
\texttt{README.md} & краткая инструкция по запуску \\
\bottomrule
\end{tabular}
\end{table}

Запуск тестов:
\begin{lstlisting}
npm test
\end{lstlisting}

Запуск локального демо:
\begin{lstlisting}
npm run serve
\end{lstlisting}

После запуска локального сервера страница открывается по адресу:
\begin{lstlisting}
http://localhost:8000
\end{lstlisting}

Визуальная часть состоит из двух основных областей. Слева расположены элементы управления
параметрами, справа -- круговая цветная картина интерференции и график \(I(r)\). При изменении
ползунков расчет выполняется заново, поэтому можно сразу наблюдать, как меняется расстояние между
кольцами, цвет картины и контраст.

\section{Проверка корректности}

Корректность модели проверялась несколькими способами. Во-первых, используются автоматические тесты,
которые проверяют базовые физические свойства:
\begin{itemize}
  \item темный центр в отраженном свете при \(t_0=0\);
  \item комплементарность отраженной и проходящей картины;
  \item максимум первого светлого кольца;
  \item формулу первого ненулевого темного кольца \(r^2\approx\lambda R\);
  \item совпадение нулевой ширины спектра с монохроматическим расчетом;
  \item нормировку гауссовых спектральных весов;
  \item параксиальный предел для сагитты \(R-\sqrt{R^2-r^2}\);
  \item восстановление длины волны по лабораторной формуле через диаметры колец.
\end{itemize}

Результат запуска тестов:
\begin{lstlisting}
tests 8
pass 8
fail 0
\end{lstlisting}

Во-вторых, для проверки асимптотики сравнивались точные радиусы темных колец с формулой
\[
r_m\approx \sqrt{m\lambda R/n}.
\]
Для параметров \(R=1\,\mathrm{м}\), \(\lambda=589\,\mathrm{нм}\), \(n=1\) были получены следующие
значения.

\begin{table}[h]
\centering
\begin{tabular}{rrrr}
\toprule
\(m\) & \(r_{\mathrm{exact}}\), мм & \(r_{\mathrm{asymp}}\), мм & относительная ошибка, \% \\
\midrule
1  & 0.767463 & 0.767463 & \(7.363\cdot10^{-6}\) \\
2  & 1.085357 & 1.085357 & \(1.473\cdot10^{-5}\) \\
5  & 1.716100 & 1.716100 & \(3.681\cdot10^{-5}\) \\
10 & 2.426930 & 2.426932 & \(7.363\cdot10^{-5}\) \\
20 & 3.432195 & 3.432200 & \(1.473\cdot10^{-4}\) \\
40 & 4.853850 & 4.853864 & \(2.945\cdot10^{-4}\) \\
\bottomrule
\end{tabular}
\end{table}

Также была проверена стандартная лабораторная формула:
\[
\lambda=\frac{D_{n+m}^2-D_n^2}{4mR}.
\]
Диаметры брались из численной модели, а затем по ним восстанавливалась длина волны. Для натриевого
света с \(\lambda=589\,\mathrm{нм}\) расчет возвращает практически то же значение.

\begin{table}[h]
\centering
\begin{tabular}{rrrrr}
\toprule
\(n\) & \(m\) & \(D_n\), мм & \(D_{n+m}\), мм & \(\lambda_{\mathrm{calc}}\), нм \\
\midrule
4  & 8 & 3.0699 & 5.3171 & 588.998612 \\
8  & 8 & 4.3414 & 6.1397 & 588.997918 \\
12 & 8 & 5.3171 & 6.8644 & 588.997225 \\
16 & 8 & 6.1397 & 7.5196 & 588.996531 \\
\bottomrule
\end{tabular}
\end{table}

\section{Результаты моделирования}

Для базового режима были выбраны параметры:
\[
R=1.00\,\mathrm{м},\qquad
\lambda=550\,\mathrm{нм},\qquad
n=1,\qquad
t_0=0,\qquad
V=1.
\]

В отраженном свете центр картины получается темным. Радиусы нескольких темных и светлых колец:

\begin{table}[h]
\centering
\begin{tabular}{rrr}
\toprule
Номер \(m\) & темное кольцо, мм & светлое кольцо, мм \\
\midrule
0  & 0.000000 & 0.524404 \\
1  & 0.741620 & 0.908295 \\
2  & 1.048809 & 1.172604 \\
3  & 1.284523 & 1.387443 \\
5  & 1.658312 & 1.739252 \\
10 & 2.345206 & 2.403121 \\
\bottomrule
\end{tabular}
\end{table}

Для проверки влияния ширины спектра сравнивались монохроматический источник
\(\lambda_0=610\,\mathrm{нм}\) и квазимонохроматический источник с той же серединой спектра, но с
\(\Delta\lambda=120\,\mathrm{нм}\). Значения нормированной интенсивности в нескольких точках:

\begin{table}[h]
\centering
\begin{tabular}{rrr}
\toprule
\(r\), мм & монохроматический свет & \(\Delta\lambda=120\,\mathrm{нм}\) \\
\midrule
0.00 & 0.0000 & 0.0000 \\
0.25 & 0.1001 & 0.1021 \\
0.50 & 0.9219 & 0.9163 \\
0.75 & 0.0587 & 0.1032 \\
1.00 & 0.8203 & 0.7151 \\
1.50 & 0.6884 & 0.5614 \\
2.00 & 0.9679 & 0.5003 \\
3.00 & 0.4870 & 0.4999 \\
\bottomrule
\end{tabular}
\end{table}

Из таблицы видно, что около центра обе картины еще похожи, но на больших радиусах
квазимонохроматическая картина стремится к средней интенсивности около \(0.5\). Это соответствует
физическому смыслу: разные длины волн дают максимумы и минимумы в разных местах, и при усреднении
контраст дальних колец исчезает.

\section{Обсуждение результатов}

Модель хорошо воспроизводит основные свойства колец Ньютона. При нулевом зазоре центр в отраженном
свете темный, потому что два отраженных луча приходят в противофазе. В проходящем свете, наоборот,
центр становится светлым, так как картина является комплементарной.

Увеличение радиуса кривизны линзы приводит к увеличению радиусов колец. Это прямо следует из
формулы \(r_m^2\approx m\lambda R/n\): при большем \(R\) пленка утолщается медленнее, поэтому для
набора нужной разности хода нужно большее расстояние от центра. Увеличение длины волны действует
похожим образом: кольца становятся шире.

Ненулевой начальный зазор \(t_0\) сдвигает фазу всей картины. Поэтому центральная точка уже не
обязательно является минимумом. Это соответствует реальной ситуации, когда между линзой и пластиной
может оставаться пыль, микроскопический зазор или деформация поверхности.

Наиболее заметное отличие квазимонохроматического света от монохроматического состоит в падении
контраста. Вблизи центра разность хода мала, и разные длины волн еще дают близкие фазы. На больших
радиусах разность хода растет, максимумы для разных длин волн расходятся, и итоговая интенсивность
становится почти равной среднему уровню. Именно поэтому при широком спектре видны только несколько
центральных колец.

\section{Выводы}

В работе построена интерактивная модель колец Ньютона, основанная на двухлучевой интерференции в
тонкой пленке. Модель учитывает точную геометрию сферической линзы, показатель преломления пленки,
начальный зазор, видность интерференции и конечную ширину спектра.

Полученные результаты согласуются с аналитическими формулами для радиусов темных и светлых колец.
Проверка по асимптотике показывает, что точная геометрическая формула переходит в стандартное
соотношение \(r_m^2\approx m\lambda R/n\) при \(r\ll R\). Дополнительная проверка через диаметры
колец воспроизводит лабораторную формулу для определения длины волны.

Интерактивная визуализация позволяет наглядно увидеть, как меняются кольца при изменении радиуса
линзы, длины волны, ширины спектра и начального зазора. Особенно хорошо виден эффект конечной
ширины спектра: дальние кольца теряют контраст, а интенсивность стремится к среднему уровню.

Таким образом, программа решает поставленную задачу компьютерного моделирования и одновременно
служит удобным инструментом для качественного понимания физики колец Ньютона.

\appendix
\section{Ключевые фрагменты программы}

Ниже приведены основные фрагменты реализации. Полный код проекта находится в файлах
\texttt{src/physics.js}, \texttt{src/app.js} и \texttt{tests/physics.test.js}.

\subsection{Толщина тонкой пленки}
\begin{lstlisting}[language=JavaScript]
function filmThicknessM(radialCoordinateM, curvatureRadiusM, gapNm) {
  var r = finiteNumber(radialCoordinateM, "radialCoordinateM");
  var radius = finiteNumber(curvatureRadiusM, "curvatureRadiusM");
  var gapM = finiteNumber(gapNm || 0, "gapNm") * 1e-9;

  var root = Math.sqrt(Math.max(0, radius * radius - r * r));
  var sagittaM = r === 0 ? 0 : (r * r) / (radius + root);
  return gapM + sagittaM;
}
\end{lstlisting}

\subsection{Интенсивность интерференции}
\begin{lstlisting}[language=JavaScript]
function interferenceAtThickness(thicknessM, wavelengthNm, options) {
  var settings = options || {};
  var visibility = clamp(
    settings.visibility === undefined ? 1 : finiteNumber(settings.visibility, "visibility"),
    0,
    1
  );
  var phase = phaseRadians(thicknessM, wavelengthNm, settings.refractiveIndex || 1);
  var cosine = Math.cos(phase);

  if (settings.mode === "reflected") {
    return 0.5 * (1 - visibility * cosine);
  }
  if (settings.mode === "transmitted") {
    return 0.5 * (1 + visibility * cosine);
  }
}
\end{lstlisting}

\subsection{Спектральное усреднение}
\begin{lstlisting}[language=JavaScript]
function spectralIntensityAtRadius(radiusM, spectrum, options) {
  var settings = options || {};
  var samples = spectrum && spectrum.length ?
    spectrum :
    spectrumSamples(settings.wavelengthNm, 0, 1);
  var thickness = filmThicknessM(
    radiusM,
    settings.curvatureRadiusM,
    settings.gapNm || 0
  );
  var intensity = 0;

  for (var i = 0; i < samples.length; i += 1) {
    intensity += samples[i].weight *
      interferenceAtThickness(thickness, samples[i].wavelengthNm, settings);
  }
  return intensity;
}
\end{lstlisting}

\subsection{Пример автоматического теста}
\begin{lstlisting}[language=JavaScript]
test("first noncentral reflected dark ring follows the Newton formula", () => {
  const params = {
    mode: "reflected",
    curvatureRadiusM: 1,
    wavelengthNm: 550,
    refractiveIndex: 1,
    gapNm: 0
  };
  const radiusM = physics.ringRadiusM(1, { ...params, kind: "dark" });
  const paraxialRadiusM = Math.sqrt(
    params.curvatureRadiusM * params.wavelengthNm * 1e-9
  );
  const intensity = physics.interferenceIntensity(radiusM, params);

  assert.ok(Math.abs(radiusM - paraxialRadiusM) / paraxialRadiusM < 1e-7);
  assert.ok(intensity < 1e-9);
});
\end{lstlisting}

\end{document}
